% BHU PYQ -- Professional Exam Template
\documentclass[11pt,a4paper]{article}

\usepackage[a4paper,top=2.5cm,bottom=2.5cm,left=2.8cm,right=2.8cm,headheight=14pt]{geometry}
\usepackage{amsmath,amssymb}
\usepackage[T1]{fontenc}
\usepackage[utf8]{inputenc}
\usepackage{lmodern}
\usepackage{microtype}
\usepackage{fancyhdr}
\usepackage{enumitem}
\usepackage{hyperref}
\usepackage{lastpage}
\usepackage{parskip}

\hypersetup{colorlinks=true,urlcolor=black,linkcolor=black,
  pdftitle={B.Sc. (Hons.) Botany Sem-V BOB-501 -- BHU PYQ}}

\pagestyle{fancy}
\fancyhf{}
\renewcommand{\headrulewidth}{0.4pt}
\renewcommand{\footrulewidth}{0.4pt}
\fancyhead[L]{\small\textit{Banaras Hindu University}}
\fancyhead[R]{\small\textit{Botany Paper: BOB-501}}
\fancyfoot[C]{\small Page \thepage\ of \pageref{LastPage}}

\newlist{parts}{enumerate}{2}
\setlist[parts,1]{label=(\alph*),leftmargin=*,itemsep=4pt,topsep=3pt}
\setlist[parts,2]{label=(\roman*),leftmargin=*,itemsep=2pt,topsep=2pt}
\newcommand{\pts}[1]{\hfill\textit{[#1]}}

\begin{document}

\begin{center}
  {\large\bfseries BANARAS HINDU UNIVERSITY}\\[2pt]
  {\normalsize B.Sc. (Hons.) Semester V Examination, 2013-14}\\[6pt]
  \rule{0.8\linewidth}{0.5pt}\\[6pt]
  {\large\bfseries Botany}\\[4pt]
  {\large\bfseries Paper: BOB-501 --- Comparative Studies of Cryptogams}\\[4pt]
  {\normalsize Time: Three hours \hfill Max. Marks: 70}
\end{center}

\rule{\linewidth}{0.4pt}

\smallskip
\noindent\textit{Note: Attempt five questions in all, selecting one question each from Section-A, B, C, and D. Section-E is compulsory. The figures in the right-hand margin indicate marks.}

\bigskip
\begin{center}
  \textbf{SECTION -- A}
\end{center}

\medskip\noindent\textbf{Question 1.} Give an illustrated account of sexual reproduction and post-fertilization development in Florideophyceae.\pts{15}

\medskip\noindent\textbf{Question 2.} Write notes on any two of the following:\pts{$7.5 \times 2 = 15$}
\begin{parts}
  \item Pigmentation in algae
  \item Life history types in Phaeophyceae
  \item Volvocine line of evolution
  \item Methods of asexual reproduction in Cyanophyceae
\end{parts}

\bigskip
\begin{center}
  \textbf{SECTION -- B}
\end{center}

\medskip\noindent\textbf{Question 3.} Give an account of Ainsworth's system of classification of fungi. Also mention the salient features of different classes.\pts{15}

\medskip\noindent\textbf{Question 4.} Write notes on any two of the following:\pts{$7.5 \times 2 = 15$}
\begin{parts}
  \item Types of fruiting bodies in Ascomycotina
  \item Types of fungal spores
  \item Economic importance of fungi
  \item Types of life-cycles in fungi
\end{parts}

\bigskip
\begin{center}
  \textbf{SECTION -- C}
\end{center}

\medskip\noindent\textbf{Question 5.} Give an illustrated account of evolution of sporophytes in bryophytes by progressive sterilization of sporogenous tissues.\pts{15}

\medskip\noindent\textbf{Question 6.} Write critical notes on any two of the following:\pts{$7.5 \times 2 = 15$}
\begin{parts}
  \item Classification of bryophytes (Smith's System)
  \item Affinities of bryophytes with algae
  \item Vegetative and asexual reproduction in bryophytes
  \item Range of gametophyte structure in thalloid bryophytes
\end{parts}

\bigskip
\begin{center}
  \textbf{SECTION -- D}
\end{center}

\medskip\noindent\textbf{Question 7.} Describe the evolution of stelar system in pteridophytes.\pts{15}

\medskip\noindent\textbf{Question 8.} Write comprehensive notes on any two of the following:\pts{$7.5 \times 2 = 15$}
\begin{parts}
  \item \textit{Marsilea} sporocarp
  \item \textit{Psilotum} sporophyte
  \item Types of sori in ferns
  \item Morphology of \textit{Ophioglossum} sporophyte
\end{parts}

\pagebreak
\begin{center}
  \textbf{SECTION -- E}
\end{center}

\medskip\noindent\textbf{Question 9.} Differentiate between and answer the following:

\smallskip
\noindent(a) Differentiate between any five of the following:\pts{$1 \times 5 = 5$}
\begin{parts}
  \item Unilocular and plurilocular sporangia
  \item Trichoblast and gonimoblast
  \item Smut and rust fungi
  \item Pseudoelaters and elaters
  \item Leptosporangiate and eusporangiate ferns
\end{parts}

\smallskip
\noindent(b) Give very brief answers for the following (in not more than 100 words each):\pts{$1 \times 5 = 5$}
\begin{parts}
  \item Name the food reserves present in Cyanophyceae.
  \item Function of heterocyst.
  \item Name two edible fungi.
  \item Write the name of a pteridophyte which can be used as biofertilizer.
  \item Write the ecological importance of \textit{Sphagnum}.
\end{parts}

\vfill
\rule{\linewidth}{0.4pt}
\begin{center}\small
\href{https://bhu-exam.vercel.app/}{Click Here}\\[4pt]\href{https://me-aryan.vercel.app/}{Click Here}\\[4pt]\href{https://quashan-me.vercel.app/}{Click Here}
\end{center}

\end{document}
